\chapter{第1章：模群、基本域与模形式 习题}
\difficultykey

\section{基础题 \easy}

\begin{problem}{矩阵与生成元}
\easy
设
\[
S=\begin{pmatrix}0&-1\\1&0\end{pmatrix},\qquad
T=\begin{pmatrix}1&1\\0&1\end{pmatrix}.
\]
计算 $S^2$、$T^3$ 与 $(ST)^3$。
\end{problem}
\begin{solution}
直接计算得
\[
S^2=\begin{pmatrix}0&-1\\1&0\end{pmatrix}^2=
\begin{pmatrix}-1&0\\0&-1\end{pmatrix}=-I.
\]
又因为
\[
T^3=\begin{pmatrix}1&1\\0&1\end{pmatrix}^3=
\begin{pmatrix}1&3\\0&1\end{pmatrix},
\]
所以 $T^3$ 对应于平移 $\tau\mapsto \tau+3$。

再算
\[
ST=\begin{pmatrix}0&-1\\1&1\end{pmatrix},
\]
于是
\[
(ST)^2=\begin{pmatrix}-1&-1\\1&0\end{pmatrix},\qquad
(ST)^3=\begin{pmatrix}-1&0\\0&-1\end{pmatrix}=-I.
\]
因此在 $\SL_2(\Z)$ 中有关系 $S^2=(ST)^3=-I$。
\end{solution}

\begin{problem}{群作用检验}
\easy
设
\[
\gamma_1=\begin{pmatrix}a&b\\c&d\end{pmatrix},\qquad
\gamma_2=\begin{pmatrix}e&f\\g&h\end{pmatrix}
\]
属于 $\SL_2(\Z)$。证明分式线性变换满足
\[
(\gamma_1\gamma_2)\tau=\gamma_1(\gamma_2\tau).
\]
\end{problem}
\begin{solution}
先算右边：
\[
\gamma_1(\gamma_2\tau)=
\frac{a\dfrac{e\tau+f}{g\tau+h}+b}{c\dfrac{e\tau+f}{g\tau+h}+d}
=\frac{a(e\tau+f)+b(g\tau+h)}{c(e\tau+f)+d(g\tau+h)}.
\]
整理得
\[
\gamma_1(\gamma_2\tau)=
\frac{(ae+bg)\tau+(af+bh)}{(ce+dg)\tau+(cf+dh)}.
\]
而
\[
\gamma_1\gamma_2=
\begin{pmatrix}ae+bg&af+bh\\ce+dg&cf+dh\end{pmatrix},
\]
所以
\[
(\gamma_1\gamma_2)\tau=
\frac{(ae+bg)\tau+(af+bh)}{(ce+dg)\tau+(cf+dh)}=
\gamma_1(\gamma_2\tau).
\]
这就证明了它确实给出一个群作用。
\end{solution}

\begin{problem}{虚部公式}
\easy
证明对任意
\[
\gamma=\begin{pmatrix}a&b\\c&d\end{pmatrix}\in \SL_2(\Z),\qquad \tau\in \HH,
\]
都有
\[
\Im(\gamma\tau)=\frac{\Im \tau}{|c\tau+d|^2}.
\]
并用它计算 $\dfrac{2i+1}{i+1}$ 的虚部。
\end{problem}
\begin{solution}
设 $\tau=x+iy$，其中 $y>0$。注意到
\[
\gamma\tau=\frac{a\tau+b}{c\tau+d},\qquad
\overline{\gamma\tau}=\frac{a\bar\tau+b}{c\bar\tau+d}.
\]
于是
\[
\gamma\tau-\overline{\gamma\tau}
=\frac{(a\tau+b)(c\bar\tau+d)-(a\bar\tau+b)(c\tau+d)}{|c\tau+d|^2}.
\]
分子化简为
\[
(ad-bc)(\tau-\bar\tau)=\tau-\bar\tau=2iy,
\]
因为 $ad-bc=1$。故
\[
\gamma\tau-\overline{\gamma\tau}=\frac{2iy}{|c\tau+d|^2}.
\]
两边取虚部便得
\[
\Im(\gamma\tau)=\frac{y}{|c\tau+d|^2}=\frac{\Im\tau}{|c\tau+d|^2}.
\]

对 $\dfrac{2i+1}{i+1}$，这里 $c=d=1$，故
\[
\Im\!\left(\frac{2i+1}{i+1}\right)=\frac{1}{|i+1|^2}=\frac{1}{2}.
\]
直接化简也得到
\[
\frac{2i+1}{i+1}=\frac{(2i+1)(1-i)}{2}=\frac{3+i}{2},
\]
虚部确为 $1/2$。
\end{solution}

\begin{problem}{基本域判定}
\easy
设标准基本域
\[
\mathcal{F}=\{\tau\in \HH:\ |\Re\tau|\leq \tfrac12,\ |\tau|\geq 1\}.
\]
判断下列点是否属于 $\mathcal{F}$：
\[
i,\qquad \frac12+\frac{i}{2},\qquad 2i,\qquad \frac34+i.
\]
\end{problem}
\begin{solution}
逐个检查定义。

对 $i$，有 $\Re(i)=0$，且 $|i|=1$，所以 $i\in\mathcal{F}$。

对 $\dfrac12+\dfrac{i}{2}$，实部满足 $\left|\Re\tau\right|=\dfrac12$，但
\[
\left|\frac12+\frac{i}{2}\right|=\sqrt{\frac14+\frac14}=\frac{1}{\sqrt2}<1,
\]
所以它不在 $\mathcal{F}$ 中。

对 $2i$，有 $\Re(2i)=0$，且 $|2i|=2\ge 1$，故 $2i\in\mathcal{F}$。

对 $\dfrac34+i$，有
\[
\left|\Re\left(\frac34+i\right)\right|=\frac34>\frac12,
\]
因此它不在 $\mathcal{F}$ 中。
\end{solution}

\begin{problem}{简单归约}
\easy
把点
\[
\tau=\frac32+2i
\]
用 $T$ 的幂归约到竖条带 $|\Re\tau|\leq \tfrac12$ 中，并判断所得点是否已经落在基本域 $\mathcal{F}$ 中。
\end{problem}
\begin{solution}
因为 $T^{-1}\tau=\tau-1$，所以
\[
T^{-1}\left(\frac32+2i\right)=\frac12+2i.
\]
这个点已经满足 $|\Re\tau|\leq \tfrac12$。

再检验模长：
\[
\left|\frac12+2i\right|^2=\frac14+4=\frac{17}{4}>1,
\]
因此
\[
\frac12+2i\in \mathcal{F}.
\]
所以所求归约点就是 $\dfrac12+2i$。
\end{solution}

\begin{problem}{无穷尖点的稳定子}
\easy
证明在 $\SL_2(\Z)$ 中，尖点 $\infty$ 的稳定子为
\[
\operatorname{Stab}_{\SL_2(\Z)}(\infty)=
\left\{\pm\begin{pmatrix}1&n\\0&1\end{pmatrix}:n\in\Z\right\}.
\]
并说明在 $\PSL_2(\Z)$ 中它由 $T$ 生成。
\end{problem}
\begin{solution}
设
\[
\gamma=\begin{pmatrix}a&b\\c&d\end{pmatrix}\in \SL_2(\Z).
\]
它作用在 $\infty$ 上的结果是
\[
\gamma\infty=\frac{a}{c}
\]
（若 $c=0$ 则理解为 $\infty$）。因此 $\gamma\infty=\infty$ 当且仅当 $c=0$。

于是
\[
\gamma=\begin{pmatrix}a&b\\0&d\end{pmatrix},\qquad ad=1.
\]
因 $a,d\in\Z$，只能有 $a=d=1$ 或 $a=d=-1$。故
\[
\gamma=\pm\begin{pmatrix}1&b\\0&1\end{pmatrix}.
\]
把 $b$ 记作 $n$ 就得到
\[
\operatorname{Stab}_{\SL_2(\Z)}(\infty)=
\left\{\pm\begin{pmatrix}1&n\\0&1\end{pmatrix}:n\in\Z\right\}.
\]
在 $\PSL_2(\Z)$ 中，$I$ 与 $-I$ 相同，因此稳定子化成
\[
\left\{T^n:n\in\Z\right\},
\]
即由 $T$ 生成。
\end{solution}

\begin{problem}{周期函数与q展开}
\easy
设 $f$ 是上半平面上的全纯函数，并满足 $f(\tau+1)=f(\tau)$。证明存在穿孔单位圆盘上的全纯函数 $F$，使得
\[
f(\tau)=F(q),\qquad q=e^{2\pi i\tau}.
\]
若再假设 $f$ 在 $y\to\infty$ 时有界，说明 $F$ 可延拓到 $q=0$，从而 $f$ 有 $q$-展开。
\end{problem}
\begin{solution}
指数映射
\[
q=e^{2\pi i\tau}
\]
把上半平面送到穿孔单位圆盘 $0<|q|<1$。定义
\[
F(q)=f\!\left(\frac{1}{2\pi i}\log q\right).
\]
这里需要说明与对数支无关。若把 $\log q$ 换成 $\log q+2\pi i n$，则对应的 $\tau$ 改变为 $\tau+n$，而 $f(\tau+1)=f(\tau)$ 保证函数值不变，因此 $F$ 定义良好。

由于 $f$ 全纯而指数映射局部双全纯，故 $F$ 在穿孔圆盘上全纯，并满足
\[
f(\tau)=F(e^{2\pi i\tau}).
\]

若 $f$ 在 $y\to\infty$ 时有界，那么当 $q=e^{2\pi i\tau}\to 0$ 时，$F(q)=f(\tau)$ 也有界。由可去奇点定理，$F$ 在 $q=0$ 处可延拓为全纯函数。因此
\[
F(q)=\sum_{n=0}^{\infty} a_n q^n,
\]
于是
\[
f(\tau)=\sum_{n=0}^{\infty} a_n e^{2\pi i n\tau}=
\sum_{n=0}^{\infty} a_n q^n.
\]
这就是 $q$-展开。
\end{solution}

\begin{problem}{尖点形式判断}
\easy
设两个函数在尖点 $\infty$ 的 $q$-展开分别为
\[
f(\tau)=1+3q+5q^2+\cdots,
\qquad
 g(\tau)=q-2q^2+4q^3+\cdots.
\]
根据定义判断哪一个可能是尖点形式，并说明理由。
\end{problem}
\begin{solution}
尖点形式的定义是在每个尖点处全纯，并且在尖点 $\infty$ 的 $q$-展开常数项为 $0$。

对 $f$，常数项是 $1$，所以 $f$ 不是尖点形式。

对 $g$，常数项是 $0$，因此从尖点 $\infty$ 这一项看，$g$ 满足尖点形式的必要条件。若它还满足相应的模变换律并在其他尖点也全纯，那么它就可以是尖点形式。

所以在这两个例子里，只有 $g$ 可能是尖点形式。
\end{solution}

\begin{problem}{复环面的简单比较}
\easy
比较下列复环面：
\[
\C/(\Z+\Z i),\qquad \C/(2\Z+2\Z i),\qquad \C/(\Z+2\Z i).
\]
判断前两个是否同构，以及第一个与第三个是否同构。
\end{problem}
\begin{solution}
记
\[
\Lambda_1=\Z+\Z i,
\qquad
\Lambda_2=2\Z+2\Z i=2\Lambda_1.
\]
因为 $\Lambda_2=2\Lambda_1$，乘以 $2$ 的映射
\[
\C/\Lambda_1\longrightarrow \C/\Lambda_2,
\qquad z\bmod \Lambda_1\mapsto 2z\bmod \Lambda_2
\]
给出复环面同构。所以前两个同构。

再看第三个格点
\[
\Lambda_3=\Z+2\Z i.
\]
若 $\C/\Lambda_1$ 与 $\C/\Lambda_3$ 同构，则应有某个 $\alpha\in\C^{\times}$ 使得
\[
\Lambda_3=\alpha\Lambda_1.
\]
这等价于参数 $2i$ 与 $i$ 在 $\SL_2(\Z)$ 作用下等价。若存在
\[
\frac{ai+b}{ci+d}=2i,
\qquad
\begin{pmatrix}a&b\\c&d\end{pmatrix}\in\SL_2(\Z),
\]
则交叉相乘得到
\[
ai+b=2i(ci+d)=-2c+2di.
\]
比较实部与虚部可得
\[
b=-2c,\qquad a=2d.
\]
于是
\[
ad-bc=2d^2+2c^2=2(c^2+d^2),
\]
不可能等于 $1$。矛盾。

因此第一个与第三个不同构。
\end{solution}

\section{进阶题 \medium}

\begin{problem}{基本域归约定理}
\medium
证明：对任意 $\tau\in\HH$，都存在 $\gamma\in\SL_2(\Z)$ 使得 $\gamma\tau\in\mathcal{F}$。
\end{problem}
\begin{solution}
先在 $\SL_2(\Z)\tau$ 这个轨道里选择虚部最大的点，再用某个 $T^n$ 把它平移到竖条带
\[
-\frac12\leq \Re\tau\leq \frac12
\]
中。记得到的点仍为 $\tau_0$。

我们断言 $|\tau_0|\geq 1$。若不然，$|\tau_0|<1$，则作用 $S$ 得到
\[
S\tau_0=-\frac1{\tau_0}.
\]
由虚部公式，
\[
\Im(S\tau_0)=\frac{\Im\tau_0}{|\tau_0|^2}>\Im\tau_0,
\]
这与 $\tau_0$ 在轨道中虚部最大矛盾。

因此 $\tau_0$ 同时满足
\[
|\Re\tau_0|\leq \frac12,
\qquad
|\tau_0|\geq 1,
\]
即 $\tau_0\in\mathcal{F}$。故任意轨道都与基本域相交。
\end{solution}

\begin{problem}{边界识别}
\medium
说明标准基本域的两类边界是如何被识别的：
\[
\Re\tau=-\frac12 \text{ 与 } \Re\tau=\frac12,
\qquad |\tau|=1.
\]
\end{problem}
\begin{solution}
先看两条竖边。若
\[
\tau=-\frac12+iy,
\qquad y>\frac{\sqrt3}{2},
\]
则
\[
T\tau=\tau+1=\frac12+iy.
\]
所以左边界与右边界由 $T$ 配对。

再看圆弧边界。若 $|\tau|=1$，则
\[
S\tau=-\frac1\tau=-\bar\tau.
\]
因为 $|\tau|=1$，所以 $1/\tau=\bar\tau$。因此 $S$ 把单位圆弧上的点反射到对称点。特别地，圆弧左半边与右半边由 $S$ 配对。

所以标准基本域的边界识别由 $T$ 与 $S$ 两个生成元完成。这也解释了为什么把边界按这些识别粘起来后会得到商空间 $\SL_2(\Z)\backslash \HH$。
\end{solution}

\begin{problem}{椭圆点稳定子}
\medium
设
\[
\rho=e^{2\pi i/3}=\frac{-1+\sqrt3\,i}{2}.
\]
求 $i$ 与 $\rho$ 在 $\SL_2(\Z)$ 中的稳定子。
\end{problem}
\begin{solution}
先看 $i$。若
\[
\gamma=\begin{pmatrix}a&b\\c&d\end{pmatrix}
\]
满足 $\gamma i=i$，则
\[
ai+b=i(ci+d)=-c+di.
\]
比较实部与虚部得
\[
b=-c,\qquad a=d.
\]
再由行列式条件
\[
ad-bc=a^2+c^2=1,
\]
只能有 $(a,c)=(\pm1,0)$ 或 $(0,\pm1)$。故稳定子为
\[
\{\pm I,\pm S\}.
\]

再看 $\rho$。条件 $\gamma\rho=\rho$ 等价于
\[
a\rho+b=\rho(c\rho+d).
\]
利用关系 $\rho^2+\rho+1=0$，把右边化为
\[
\rho(c\rho+d)=c\rho^2+d\rho=-c+(d-c)\rho.
\]
因此
\[
b=-c,\qquad a=d-c.
\]
代回行列式条件得
\[
a(a+c)+c^2=1,
\]
等价于
\[
a^2+ac+c^2=1.
\]
整数解只有
\[
(a,c)=(1,0),\ (0,1),\ (-1,1),\ (-1,0),\ (0,-1),\ (1,-1).
\]
相应地得到六个矩阵
\[
\pm I,\ \pm ST,\ \pm (ST)^2.
\]
所以 $\rho$ 的稳定子为
\[
\{\pm I,\pm ST,\pm (ST)^2\}.
\]
在 $\PSL_2(\Z)$ 中，它分别对应阶 $2$ 与阶 $3$ 的稳定子。
\end{solution}

\begin{problem}{奇权消失}
\medium
证明：若 $k$ 为奇数，则
\[
M_k(\SL_2(\Z))=\{0\}.
\]
\end{problem}
\begin{solution}
设 $f\in M_k(\SL_2(\Z))$。因为
\[
-I=\begin{pmatrix}-1&0\\0&-1\end{pmatrix}\in \SL_2(\Z),
\]
而 $-I$ 在上半平面上的作用是恒等，所以模变换律给出
\[
f(\tau)=f|_k(-I)(\tau)=(-1)^k f(\tau).
\]
若 $k$ 为奇数，则 $(-1)^k=-1$，于是
\[
f(\tau)=-f(\tau).
\]
故对所有 $\tau\in\HH$ 都有 $f(\tau)=0$。因此
\[
M_k(\SL_2(\Z))=\{0\}.
\]
\end{solution}

\begin{problem}{指标计算}
\medium
利用公式
\[
[\SL_2(\Z):\Gamma_0(N)]=N\prod_{p\mid N}\left(1+\frac1p\right)
\]
计算 $N=2,3,4,6,11$ 时的指标。
\end{problem}
\begin{solution}
逐个代入即可。

当 $N=2$ 时，
\[
[\SL_2(\Z):\Gamma_0(2)]=2\left(1+\frac12\right)=3.
\]

当 $N=3$ 时，
\[
[\SL_2(\Z):\Gamma_0(3)]=3\left(1+\frac13\right)=4.
\]

当 $N=4$ 时，素因子只有 $2$，故
\[
[\SL_2(\Z):\Gamma_0(4)]=4\left(1+\frac12\right)=6.
\]

当 $N=6$ 时，素因子为 $2,3$，故
\[
[\SL_2(\Z):\Gamma_0(6)]=6\left(1+\frac12\right)\left(1+\frac13\right)=6\cdot\frac32\cdot\frac43=12.
\]

当 $N=11$ 时，
\[
[\SL_2(\Z):\Gamma_0(11)]=11\left(1+\frac1{11}\right)=12.
\]
\end{solution}

\begin{problem}{Gamma0二的尖点}
\medium
证明 $\Gamma_0(2)$ 只有两个尖点，可取代表为 $\infty$ 与 $0$；并求它们的宽度。
\end{problem}
\begin{solution}
任取尖点 $a/c\in \mathbb{P}^1(\Q)$，其中 $(a,c)=1$。

若 $c$ 为偶数，则存在整数 $b,d$ 使得
\[
ad-bc=1.
\]
于是矩阵
\[
\gamma=\begin{pmatrix}a&b\\c&d\end{pmatrix}
\]
属于 $\Gamma_0(2)$，并且
\[
\gamma\infty=\frac{a}{c}.
\]
所以这类尖点都与 $\infty$ 等价。

若 $c$ 为奇数，则 $a$ 也必为奇数。因为 $(2a,c)=1$，可解方程
\[
xc-2ay=1.
\]
于是矩阵
\[
\gamma=\begin{pmatrix}x&a\\2y&c\end{pmatrix}
\]
满足行列式为 $1$，且下左角被 $2$ 整除，所以 $\gamma\in\Gamma_0(2)$。又
\[
\gamma 0=\frac{a}{c}.
\]
因此这类尖点都与 $0$ 等价。

所以 $\Gamma_0(2)$ 只有两个尖点：$\infty$ 与 $0$。

宽度方面，$\infty$ 的宽度是使 $T^n\in\Gamma_0(2)$ 的最小正整数。由于
\[
T^n=\begin{pmatrix}1&n\\0&1\end{pmatrix}
\]
总在 $\Gamma_0(2)$ 中，所以宽度为 $1$。

对 $0$，用 $S$ 把它送到 $\infty$。有
\[
ST^nS^{-1}=\begin{pmatrix}1&0\\-n&1\end{pmatrix}.
\]
它属于 $\Gamma_0(2)$ 当且仅当 $n$ 为偶数。因此 $0$ 的宽度为 $2$。
\end{solution}

\begin{problem}{Gamma0三的尖点}
\medium
证明 $\Gamma_0(3)$ 只有两个尖点，可取代表为 $\infty$ 与 $0$；并求它们的宽度。
\end{problem}
\begin{solution}
论证与 $N=2$ 时相同。

任取尖点 $a/c$，其中 $(a,c)=1$。

若 $3\mid c$，则可取
\[
\gamma=\begin{pmatrix}a&b\\c&d\end{pmatrix}\in \Gamma_0(3)
\]
使得 $ad-bc=1$，于是 $\gamma\infty=a/c$，故该尖点与 $\infty$ 等价。

若 $3\nmid c$，则 $(3a,c)=1$，可解
\[
xc-3ay=1.
\]
于是
\[
\gamma=\begin{pmatrix}x&a\\3y&c\end{pmatrix}\in\Gamma_0(3),
\qquad
\gamma 0=\frac{a}{c}.
\]
所以该尖点与 $0$ 等价。

因此 $\Gamma_0(3)$ 只有两个尖点：$\infty$ 与 $0$。

宽度方面，$\infty$ 仍由 $T$ 生成，所以宽度为 $1$。

对 $0$，同样考察
\[
ST^nS^{-1}=\begin{pmatrix}1&0\\-n&1\end{pmatrix}.
\]
它落在 $\Gamma_0(3)$ 中当且仅当 $3\mid n$，所以 $0$ 的宽度为 $3$。
\end{solution}

\begin{problem}{格点与模群作用}
\medium
设
\[
\Lambda_\tau=\Z+\Z\tau,
\qquad
\gamma=\begin{pmatrix}a&b\\c&d\end{pmatrix}\in\SL_2(\Z).
\]
证明
\[
\Lambda_{\gamma\tau}=(c\tau+d)^{-1}\Lambda_\tau.
\]
从而说明 $\C/\Lambda_\tau$ 与 $\C/\Lambda_{\gamma\tau}$ 同构。
\end{problem}
\begin{solution}
由定义
\[
\gamma\tau=\frac{a\tau+b}{c\tau+d}.
\]
于是
\[
\Lambda_{\gamma\tau}=\Z+\Z\frac{a\tau+b}{c\tau+d}.
\]
把右边提出公因子 $(c\tau+d)^{-1}$，得
\[
\Lambda_{\gamma\tau}
=(c\tau+d)^{-1}\bigl((c\tau+d)\Z+(a\tau+b)\Z\bigr).
\]
而
\[
(c\tau+d)\Z+(a\tau+b)\Z
\]
由 $1$ 与 $\tau$ 的整系数组合生成，因为
\[
a\tau+b=a\tau+b\cdot 1,
\qquad
c\tau+d=c\tau+d\cdot 1.
\]
反过来，由于
\[
\begin{pmatrix}a&b\\c&d\end{pmatrix}
\in\SL_2(\Z)
\]
可逆，$1$ 与 $\tau$ 也能写成这两个数的整系数组合。因此
\[
(c\tau+d)\Z+(a\tau+b)\Z=\Z+\Z\tau=\Lambda_\tau.
\]
故
\[
\Lambda_{\gamma\tau}=(c\tau+d)^{-1}\Lambda_\tau.
\]
于是乘以 $(c\tau+d)$ 给出格点之间的比例变换，所以
\[
\C/\Lambda_\tau\cong \C/\Lambda_{\gamma\tau}.
\]
\end{solution}

\begin{problem}{椭圆曲线的j值}
\medium
对椭圆曲线
\[
E_1: y^2=x^3-x,
\qquad
E_2: y^2=x^3+1,
\]
计算它们的判别式与 $j$-不变量。
\end{problem}
\begin{solution}
对 Weierstrass 方程
\[
y^2=x^3+Ax+B,
\]
有公式
\[
\Delta_E=-16(4A^3+27B^2),
\qquad
j(E)=1728\,\frac{4A^3}{4A^3+27B^2}.
\]

对 $E_1$，有 $A=-1$，$B=0$，故
\[
4A^3+27B^2=4(-1)^3=-4,
\qquad
\Delta_{E_1}=-16(-4)=64.
\]
于是
\[
j(E_1)=1728\cdot\frac{-4}{-4}=1728.
\]

对 $E_2$，有 $A=0$，$B=1$，故
\[
4A^3+27B^2=27,
\qquad
\Delta_{E_2}=-16\cdot 27=-432.
\]
于是
\[
j(E_2)=1728\cdot\frac{0}{27}=0.
\]
因此
\[
\Delta_{E_1}=64,\ j(E_1)=1728;
\qquad
\Delta_{E_2}=-432,\ j(E_2)=0.
\]
\end{solution}

\section{挑战题 \hard}

\begin{problem}{双曲度量不变性}
\hard
证明双曲度量
\[
ds^2=\frac{dx^2+dy^2}{y^2}
\]
在 $\SL_2(\Z)$ 的分式线性变换下不变。
\end{problem}
\begin{solution}
把 $\tau=x+iy$ 看成复变量，双曲度量可写成
\[
ds^2=\frac{|d\tau|^2}{(\Im\tau)^2}.
\]
设
\[
\gamma\tau=\frac{a\tau+b}{c\tau+d}.
\]
对 $\tau$ 求导得
\[
d(\gamma\tau)=\frac{(ad-bc)\,d\tau}{(c\tau+d)^2}=\frac{d\tau}{(c\tau+d)^2},
\]
因为 $ad-bc=1$。于是
\[
|d(\gamma\tau)|^2=\frac{|d\tau|^2}{|c\tau+d|^4}.
\]
另一方面，由虚部公式
\[
\Im(\gamma\tau)=\frac{\Im\tau}{|c\tau+d|^2}.
\]
所以
\[
\frac{|d(\gamma\tau)|^2}{\Im(\gamma\tau)^2}
=\frac{|d\tau|^2/|c\tau+d|^4}{(\Im\tau)^2/|c\tau+d|^4}
=\frac{|d\tau|^2}{(\Im\tau)^2}.
\]
因此双曲度量在 $\gamma$ 下保持不变，即
\[
\frac{dx'^2+dy'^2}{y'^2}=\frac{dx^2+dy^2}{y^2}.
\]
\end{solution}

\begin{problem}{内部点的唯一性}
\hard
证明：若 $\tau$ 与 $\gamma\tau$ 都位于基本域 $\mathcal{F}$ 的内部，则 $\gamma=\pm I$。
\end{problem}
\begin{solution}
设
\[
\gamma=\begin{pmatrix}a&b\\c&d\end{pmatrix}\in\SL_2(\Z),
\]
并假设 $\tau,\gamma\tau$ 都在 $\mathcal{F}$ 的内部。内部条件是
\[
|\Re\tau|<\frac12,
\qquad
|\tau|>1.
\]

若 $c\neq 0$，则对基本域内部的点有标准事实
\[
|c\tau+d|>1.
\]
于是由虚部公式可得
\[
\Im(\gamma\tau)=\frac{\Im\tau}{|c\tau+d|^2}<\Im\tau.
\]
但把同样论证应用到 $\gamma^{-1}$ 与点 $\gamma\tau$，又得到
\[
\Im\tau<\Im(\gamma\tau),
\]
矛盾。因此必有 $c=0$。

于是
\[
\gamma=\begin{pmatrix}a&b\\0&d\end{pmatrix},
\qquad ad=1,
\]
所以 $a=d=\pm1$。模去整体符号后，可写成 $T^n$。也就是说，
\[
\gamma=\pm \begin{pmatrix}1&n\\0&1\end{pmatrix}.
\]
若 $n\neq 0$，则 $\gamma\tau=\tau+n$ 的实部与 $\tau$ 的实部相差整数，不可能同时都落在开区间 $(-1/2,1/2)$ 中。故只能有 $n=0$。

因此
\[
\gamma=\pm I.
\]
\end{solution}

\begin{problem}{尖点等于有理射影直线}
\hard
证明模群的尖点集合正好是
\[
\mathbb{P}^1(\Q)=\Q\cup\{\infty\}.
\]
\end{problem}
\begin{solution}
按定义，尖点是 $\SL_2(\Z)$ 对 $\infty$ 的像的全体。

若
\[
\gamma=\begin{pmatrix}a&b\\c&d\end{pmatrix}\in\SL_2(\Z),
\]
则
\[
\gamma\infty=
\begin{cases}
\infty,& c=0,\\[4pt]
a/c,& c\neq 0.
\end{cases}
\]
因此每个尖点都在 $\Q\cup\{\infty\}$ 中。

反过来，任取有理数 $a/c\in\Q$，其中 $(a,c)=1$。由 Bézout 等式可取整数 $b,d$ 使
\[
ad-bc=1.
\]
于是
\[
\gamma=\begin{pmatrix}a&b\\c&d\end{pmatrix}\in\SL_2(\Z),
\qquad
\gamma\infty=a/c.
\]
而 $\infty$ 本身当然也是尖点。

故尖点集合恰为
\[
\mathbb{P}^1(\Q)=\Q\cup\{\infty\}.
\]
\end{solution}

\begin{problem}{非平凡稳定子的点}
\hard
证明：若 $\tau\in\HH$ 的稳定子在 $\PSL_2(\Z)$ 中非平凡，则 $\tau$ 与 $i$ 或 $\rho=e^{2\pi i/3}$ 等价。
\end{problem}
\begin{solution}
设 $\tau$ 被某个非平凡元素 $\gamma\in\PSL_2(\Z)$ 固定。取 $\eta\in\PSL_2(\Z)$ 使
\[
\tau_0=\eta\tau\in\mathcal{F}.
\]
则
\[
\eta\gamma\eta^{-1}\tau_0=\tau_0,
\]
所以 $\tau_0$ 也有非平凡稳定子。

若 $\tau_0$ 在 $\mathcal{F}$ 内部，则由上一题知只有恒等元能固定它，矛盾。因此 $\tau_0$ 必在边界上。

再看边界各部分：

竖边上的开线段只由 $T$ 与相邻边配对，但 $T\tau_0=\tau_0+1\neq\tau_0$，所以这些点没有非平凡稳定子。

单位圆弧上的点由 $S$ 配对。若 $S\tau_0=\tau_0$，则
\[
-\frac1{\tau_0}=\tau_0,
\]
故 $\tau_0^2=-1$，于是 $\tau_0=i$。

端点处，$\rho$ 满足 $(ST)\rho=\rho$，因此有阶 $3$ 的稳定子。另一个端点 $e^{\pi i/3}$ 与 $\rho$ 等价，所以不需要另列。

因此所有具有非平凡稳定子的点都与 $i$ 或 $\rho$ 等价。
\end{solution}

\begin{problem}{模群的表示}
\hard
说明为什么
\[
\PSL_2(\Z)\cong \Z/2\Z * \Z/3\Z.
\]
\end{problem}
\begin{solution}
由第 1 题可知 $\SL_2(\Z)$ 由 $S,T$ 生成，并满足
\[
S^2=(ST)^3=-I.
\]
在 $\PSL_2(\Z)$ 中把 $\pm I$ 视为同一个元素，于是记
\[
s=[S],\qquad u=[ST].
\]
则
\[
s^2=1,\qquad u^3=1.
\]
而且 $[T]=su^{-1}$，所以 $s,u$ 生成整个 $\PSL_2(\Z)$。

另一方面，标准基本域是一个有两个顶点角分别为 $\pi/2$ 与 $\pi/3$ 的双曲多边形，边配对由 $S$ 与 $T$ 给出。对应的唯一基本关系恰好就是顶点稳定子带来的
\[
s^2=1,\qquad u^3=1.
\]
没有额外关系。

因此 $\PSL_2(\Z)$ 由一个阶 $2$ 元与一个阶 $3$ 元自由积生成：
\[
\PSL_2(\Z)\cong \langle s,u\mid s^2=1,\ u^3=1\rangle
\cong \Z/2\Z * \Z/3\Z.
\]
这正是所求结论。
\end{solution}

\begin{problem}{Gamma0四的尖点}
\hard
证明 $\Gamma_0(4)$ 有三个尖点，可取代表为
\[
\infty,\qquad 0,\qquad \frac12,
\]
并求它们的宽度。
\end{problem}
\begin{solution}
任取尖点 $a/c$，其中 $(a,c)=1$。按 $c$ 模 $4$ 的情形分类。

若 $4\mid c$，则可取
\[
\gamma=\begin{pmatrix}a&b\\c&d\end{pmatrix}\in\Gamma_0(4)
\]
使 $ad-bc=1$，于是 $\gamma\infty=a/c$，所以该尖点与 $\infty$ 等价。

若 $c$ 为奇数，则 $(4a,c)=1$，可解
\[
xc-4ay=1.
\]
于是
\[
\gamma=\begin{pmatrix}x&a\\4y&c\end{pmatrix}\in\Gamma_0(4),
\qquad
\gamma 0=\frac{a}{c},
\]
故该尖点与 $0$ 等价。

若 $c\equiv 2\pmod 4$，写 $c=2c_1$，其中 $c_1$ 为奇数。设
\[
\sigma=\begin{pmatrix}1&0\\2&1\end{pmatrix},
\qquad \sigma\infty=\frac12.
\]
我们要找 $\gamma\in\Gamma_0(4)$ 使 $\gamma\sigma\infty=a/c$。取整数 $s,v$ 满足
\[
av-cs=1.
\]
因为 $c$ 为偶数、$a$ 为奇数，所以 $v$ 必为奇数。再令
\[
u=\frac{c-2v}{4},\qquad r=a-2s.
\]
则
\[
\gamma=\begin{pmatrix}r&s\\4u&v\end{pmatrix}\in\Gamma_0(4),
\]
并且
\[
\gamma\sigma=
\begin{pmatrix}r+2s&s\\4u+2v&v\end{pmatrix}=
\begin{pmatrix}a&s\\c&v\end{pmatrix},
\]
所以
\[
\gamma\sigma\infty=\frac{a}{c}.
\]
因此这类尖点都与 $1/2$ 等价。

故 $\Gamma_0(4)$ 的尖点恰有三个：$\infty,0,1/2$。

下面求宽度。

对 $\infty$，因为 $T\in\Gamma_0(4)$，所以宽度为 $1$。

对 $0$，仍有
\[
ST^nS^{-1}=\begin{pmatrix}1&0\\-n&1\end{pmatrix}.
\]
它属于 $\Gamma_0(4)$ 当且仅当 $4\mid n$，故 $0$ 的宽度为 $4$。

对 $1/2$，用 $\sigma$ 共轭：
\[
\sigma T^n\sigma^{-1}=
\begin{pmatrix}1-2n&n\\-4n&2n+1\end{pmatrix}.
\]
其左下角总被 $4$ 整除，所以最小正整数 $n=1$ 就可行。因此 $1/2$ 的宽度为 $1$。
\end{solution}

\begin{problem}{低权单项式}
\hard
已知 $E_4$ 是权 $4$ 模形式，$E_6$ 是权 $6$ 模形式。证明任意单项式
\[
E_4^aE_6^b
\]
都是权 $4a+6b$ 的模形式，并列出权不超过 $12$ 时由这些单项式得到的所有例子。
\end{problem}
\begin{solution}
若 $f$ 是权 $k$ 模形式，$g$ 是权 $\ell$ 模形式，则它们的乘积满足
\[
(fg)(\gamma\tau)=f(\gamma\tau)g(\gamma\tau)
=(c\tau+d)^k f(\tau)\,(c\tau+d)^\ell g(\tau)
=(c\tau+d)^{k+\ell}f(\tau)g(\tau).
\]
所以 $fg$ 是权 $k+\ell$ 模形式。反复应用便得
\[
E_4^aE_6^b\in M_{4a+6b}(\SL_2(\Z)).
\]

当权不超过 $12$ 时，可出现的非负整数解如下：

权 $0$：$1$；

权 $4$：$E_4$；

权 $6$：$E_6$；

权 $8$：$E_4^2$；

权 $10$：$E_4E_6$；

权 $12$：$E_4^3$ 与 $E_6^2$。

因此低权中最自然的候选模形式正是
\[
1,\ E_4,\ E_6,\ E_4^2,\ E_4E_6,\ E_4^3,\ E_6^2.
\]
这说明 $E_4,E_6$ 很像整个分次环的生成元；后续章节会把这一现象精确化。
\end{solution}

\begin{problem}{j不变量与复椭圆曲线}
\hard
说明为什么两个复椭圆曲线同构当且仅当它们有相同的 $j$-不变量。
\end{problem}
\begin{solution}
每条复椭圆曲线都可以写成某个复环面
\[
E\cong \C/\Lambda.
\]
若两个格点满足
\[
\Lambda'=\alpha\Lambda\qquad (\alpha\in\C^{\times}),
\]
则乘以 $\alpha$ 给出复环面同构
\[
\C/\Lambda\cong \C/\Lambda'.
\]
因此复椭圆曲线的同构类与格点的相似类对应。

另一方面，每个格点相似类都可写成
\[
\Lambda_\tau=\Z+\Z\tau,
\qquad \tau\in\HH,
\]
而且两个参数 $\tau,\tau'$ 给出同构复环面，当且仅当它们在 $\SL_2(\Z)$ 作用下等价。$j$-函数正是这个商空间上的全纯不变量：
\[
j(\gamma\tau)=j(\tau).
\]
因此它在同构的复环面上取同一值。

反过来，章节中的分类定理说明：$j$-不变量完全决定复椭圆曲线的同构类。也就是说，若
\[
j(E_1)=j(E_2),
\]
则对应的格点相似类相同，从而
\[
E_1\cong E_2.
\]
所以“同构当且仅当 $j$ 相同”。
\end{solution}

\begin{problem}{格点的标准形与唯一性}
\hard
证明任意格点 $\Lambda\subset\C$ 都与某个
\[
\Lambda_\tau=\Z+\Z\tau,
\qquad \tau\in\mathcal{F}
\]
相似；并说明当 $\tau$ 取在基本域内部时，这个表示是唯一的。
\end{problem}
\begin{solution}
任取格点 $\Lambda$ 的一组基底 $\omega_1,\omega_2$，并设
\[
\tau=\frac{\omega_2}{\omega_1}.
\]
调整基底顺序后可令 $\Im\tau>0$，于是
\[
\Lambda=\omega_1(\Z+\Z\tau)=\omega_1\Lambda_\tau.
\]
这说明每个格点都与某个 $\Lambda_\tau$ 相似。

由基本域归约定理，存在 $\gamma\in\SL_2(\Z)$ 使得 $\gamma\tau\in\mathcal{F}$。上一节已经证明
\[
\Lambda_{\gamma\tau}=(c\tau+d)^{-1}\Lambda_\tau,
\]
所以 $\Lambda$ 也与某个 $\Lambda_{\tau_0}$ 相似，其中 $\tau_0\in\mathcal{F}$。

再说明唯一性。若
\[
\Lambda_{\tau_1}\sim \Lambda_{\tau_2},
\qquad \tau_1,\tau_2\in \mathcal{F}^{\circ}
\]
都在基本域内部，则必有某个 $\gamma\in\SL_2(\Z)$ 使
\[
\tau_2=\gamma\tau_1.
\]
由“内部点的唯一性”知 $\gamma=\pm I$，于是 $\tau_2=\tau_1$。

因此，每个格点都有一个落在 $\mathcal{F}$ 中的标准参数，而在基本域内部该参数唯一。
\end{solution}
